\documentclass[sigconf]{acmart}

\usepackage{graphicx}
\usepackage{booktabs}
\usepackage{hyperref}

\AtBeginDocument{%
  \providecommand\BibTeX{{%
    Bib\TeX}}}

\setcopyright{acmlicensed}
\copyrightyear{2026}
\acmYear{2026}
\acmDOI{10.1145/XXXXXXX.XXXXXXX}
\acmConference[WSDM '26]{Proceedings of the 19th ACM International Conference on Web Search and Data Mining}{February 22--26, 2026}{Boise, ID, USA}
\acmBooktitle{Proceedings of the 19th ACM International Conference on Web Search and Data Mining (WSDM '26), February 22--26, 2026, Boise, ID, USA}
\acmISBN{978-1-4503-XXXX-X/26/02}

\begin{document}

\title{LLM-to-Map: Conversational Search and Retrieval for Urban Geospatial Data}

\author{First Author}
\email{first.author@institution.edu}
\affiliation{%
  \institution{Your Institution}
  \city{City}
  \country{Country}}

\author{Second Author}
\email{second.author@institution.edu}
\affiliation{%
  \institution{Your Institution}
  \city{City}
  \country{Country}}

\renewcommand{\shortauthors}{Author et al.}

\keywords{Natural Language Search, Geospatial Information Retrieval, Fuzzy Matching, Large Language Models, Web-based Visualization}

\begin{abstract}
We demonstrate \textbf{LLM-to-Map}, a web-based system for conversational search and retrieval of urban geospatial data. Traditional geographic information systems require exact coordinate queries and structured search interfaces. Our system applies large language models and fuzzy matching algorithms to enable fault-tolerant natural language search over urban infrastructure data. Users can search with typos (``times sqare'' → Times Square), vague descriptions (``the station near broadway''), and context-dependent queries (``zoom out'' maintains previous location). The system integrates: (1) Levenshtein distance-based fuzzy matching for location name retrieval, (2) LLM-powered query understanding with context tracking across conversational turns, (3) real-time data mining of power grid and traffic network states, (4) automated web-based visualization with data-driven camera positioning. We demonstrate the system on Manhattan's infrastructure network (8 substations, 50+ landmarks), showing how conversational search can democratize access to complex urban datasets. The system handles location ambiguity, spelling variations, and multi-turn context—key challenges in geospatial information retrieval that traditional keyword search cannot address.
\end{abstract}

\begin{CCSXML}
<ccs2012>
<concept>
<concept_id>10002951.10003317.10003347.10003350</concept_id>
<concept_desc>Information systems~Information retrieval query processing</concept_desc>
<concept_significance>500</concept_significance>
</concept>
<concept>
<concept_id>10002951.10003260.10003282</concept_id>
<concept_desc>Information systems~Geographic information systems</concept_desc>
<concept_significance>500</concept_significance>
</concept>
<concept>
<concept_id>10010147.10010257.10010293.10010294</concept_id>
<concept_desc>Computing methodologies~Natural language generation</concept_desc>
<concept_significance>300</concept_significance>
</concept>
</ccs2012>
\end{CCSXML}

\ccsdesc[500]{Information systems~Information retrieval query processing}
\ccsdesc[500]{Information systems~Geographic information systems}
\ccsdesc[300]{Computing methodologies~Natural language generation}

\maketitle

\section{Introduction}

\textbf{Problem:} Urban infrastructure datasets (power grids, transportation networks) contain rich geospatial information, but accessing this data requires technical expertise. Operators must know exact coordinates (``-73.986, 40.758''), navigate complex query interfaces, and manually filter results. This creates barriers for emergency responders, urban planners, and researchers who need rapid information retrieval.

\textbf{Challenge:} Traditional keyword search fails for geospatial queries because:
\begin{itemize}
    \item \textbf{Spelling variations:} ``Times Square'' vs ``times sqare'' vs ``time square''
    \item \textbf{Ambiguity:} ``Penn'' could mean Penn Station, Pennsylvania Avenue, or University of Pennsylvania
    \item \textbf{Context dependency:} ``zoom out'' only makes sense after a location query
    \item \textbf{Vague descriptions:} ``the station near broadway''
\end{itemize}

Recent advances in large language models~\cite{achiam2023gpt,zhou2023comprehensive} enable conversational interfaces, but applying them to geospatial search raises unique challenges: how to handle location name variations, maintain spatial context across dialogue turns, and integrate with real-time infrastructure data mining.

We present \textbf{LLM-to-Map}, a web-based system demonstrating conversational search for urban geospatial data. Our key insight: combining \textit{fuzzy string matching} for location retrieval with \textit{LLM-based query understanding} creates a robust search system that handles natural language imperfections while maintaining precision.

\subsection{Contributions}

\begin{enumerate}
    \item \textbf{Fault-tolerant geospatial search:} Levenshtein distance algorithm (≤2 edits) for location name retrieval with spelling errors
    \item \textbf{Context-aware query processing:} LLM tracks conversation state to resolve references like ``it'', ``there'', ``zoom out''
    \item \textbf{Real-time data mining integration:} Search results reflect live infrastructure states (operational/failed substations, traffic conditions)
    \item \textbf{Web-based demonstration:} Interactive system on Manhattan infrastructure (8 substations, 50+ locations)
\end{enumerate}

\section{Related Work}

\textbf{Geographic Information Retrieval.} Traditional GIS requires structured queries~\cite{longley2005geographic}. Early work on natural language interfaces used rule-based parsing~\cite{popescu2003natural}, limiting flexibility. Recent efforts apply deep learning to geographic question answering~\cite{roberts2023gpt,gao2021geographic}, but focus on static datasets without real-time data mining.

\textbf{Fuzzy String Matching.} Edit distance algorithms~\cite{navarro2001guided} are well-established for spell correction. Our contribution is applying these to \textit{geospatial entity retrieval} where location aliases and partial names create unique challenges.

\textbf{LLM for Search.} Foundation models enable conversational search~\cite{zhou2023comprehensive}, but prior work focuses on web documents~\cite{anelli2023llm}. We demonstrate LLMs for \textit{structured geospatial data retrieval} with real-time mining of infrastructure networks.

\textbf{Urban Data Mining.} Power grid~\cite{pypsa2018} and traffic~\cite{lopez2018} analysis typically use domain-specific tools. Our system mines these networks to enrich search results with live status data.

\section{System Architecture}

\subsection{Query Processing Pipeline}

\textbf{Input:} User natural language query (``show times sqare'')

\textbf{Stage 1 - Preprocessing:} Apply typo dictionary for common errors:
\begin{verbatim}
typo_corrections = {
    'confrim': 'confirm',
    'restore everyth': 'restore everything',
    'sqare': 'square',
    'staton': 'station'
}
\end{verbatim}

\textbf{Stage 2 - LLM Query Understanding:}
\begin{verbatim}
# OpenAI API call
response = openai_client.chat.completions.create(
    model="gpt-3.5-turbo",
    messages=[{
        "role": "system",
        "content": """Classify query intent:
        - location_search: Find/show a place
        - infrastructure_control: Modify system
        - data_query: Retrieve information
        - navigation: Camera/view adjustment"""
    }, {
        "role": "user",
        "content": user_query
    }]
)
intent = extract_intent(response)
\end{verbatim}

\textbf{Stage 3 - Fuzzy Location Retrieval:} For location searches, apply Levenshtein distance:
\begin{verbatim}
def fuzzy_location_search(query_text, locations_db):
    """
    Search with typo tolerance using edit distance
    Returns: (matched_location, distance_score)
    """
    min_distance = float('inf')
    best_match = None

    # Try all word combinations in query
    for word_combo in get_combinations(query_text):
        for location in locations_db:
            # Check location name and all aliases
            for name_variant in [location.name] +
                                 location.aliases:
                dist = levenshtein_distance(
                    word_combo.lower(),
                    name_variant.lower()
                )
                if dist <= 2 and dist < min_distance:
                    min_distance = dist
                    best_match = location

    return best_match, min_distance
\end{verbatim}

Implementation uses dynamic programming (O(mn) time, O(m) space):
\begin{verbatim}
def levenshtein_distance(s1, s2):
    if len(s2) == 0: return len(s1)
    previous_row = list(range(len(s2) + 1))
    for i, c1 in enumerate(s1):
        current_row = [i + 1]
        for j, c2 in enumerate(s2):
            insertions = previous_row[j+1] + 1
            deletions = current_row[j] + 1
            substitutions = previous_row[j] + (c1!=c2)
            current_row.append(
                min(insertions, deletions, substitutions)
            )
        previous_row = current_row
    return previous_row[-1]
\end{verbatim}

\textbf{Stage 4 - Data Mining:} Retrieve real-time infrastructure state:
\begin{verbatim}
# Mine power grid status
for substation in power_network.substations:
    if substation.operational:
        status = query_pypsa_state(substation)

# Mine traffic conditions
traffic_data = query_sumo_vehicles(location)
\end{verbatim}

\textbf{Stage 5 - Result Visualization:} Position web-based map camera:
\begin{verbatim}
// Mapbox GL JS - data-driven camera
map.flyTo({
    center: [location.lng, location.lat],
    zoom: calculate_zoom(location.type),
    duration: 3000
});
\end{verbatim}

\subsection{Context Tracking for Multi-Turn Search}

LLM maintains conversation state to resolve context-dependent queries:

\begin{verbatim}
conversation_context = {
    'last_location': None,
    'last_substation': None,
    'entities_mentioned': set(),
    'pending_actions': []
}

# Example:
# User: "show times square"
# → context.last_location = "Times Square"
# User: "zoom out"
# → Use context.last_location as center point
\end{verbatim}

This enables natural multi-turn searches without repeating location names.

\subsection{Location Database}

Manhattan infrastructure with geospatial metadata:

\begin{table}[h]
\centering
\caption{Location Database Schema}
\small
\begin{tabular}{lp{5cm}}
\toprule
\textbf{Field} & \textbf{Description} \\
\midrule
name & Canonical location name \\
coords & [longitude, latitude] \\
type & commercial/residential/transport/landmark \\
substation & Associated power substation \\
capacity\_mva & Power capacity (if applicable) \\
aliases & List of name variations \\
\bottomrule
\end{tabular}
\end{table}

\textbf{Example entries (from lines 214-350, ultra\_intelligent\_chatbot.py):}

\begin{itemize}
    \item \textbf{Times Square:} coords=[-73.9857, 40.7580], aliases=[``times sq'', ``time square'', ``times squere'', ``ts''], capacity=850 MVA
    \item \textbf{Penn Station:} coords=[-73.9904, 40.7505], aliases=[``penn'', ``penn st'', ``pennsylvania''], capacity=900 MVA
    \item \textbf{Grand Central:} coords=[-73.9772, 40.7527], aliases=[``grand'', ``gc'', ``grand central station''], capacity=1000 MVA
\end{itemize}

Total: 8 substations, 20+ landmarks, 10+ major streets.

\section{Demonstration Scenarios}

\subsection{Scenario 1: Fault-Tolerant Location Search}

\textbf{Query:} ``show me times sqare'' (typo: ``sqare'')

\textbf{System Flow:}
\begin{enumerate}
    \item Preprocessing: No dictionary match for ``sqare''
    \item LLM classification: intent=``location\_search'', entity=``times sqare''
    \item Fuzzy search: Levenshtein(``times sqare'', ``times square'') = 1
    \item Match found: Times Square, coords=[-73.9857, 40.7580]
    \item Data mining: Query substation status (operational)
    \item Visualization: Camera flies to location, displays power capacity
\end{enumerate}

\textbf{Result:} Successful retrieval despite spelling error. Traditional exact-match search would fail.

\subsection{Scenario 2: Context-Dependent Multi-Turn Search}

\textbf{Turn 1:} ``show penn station''
\begin{itemize}
    \item System: Displays Penn Station, stores in context
\end{itemize}

\textbf{Turn 2:} ``zoom out''
\begin{itemize}
    \item System: Uses context.last\_location = Penn Station
    \item Adjusts camera zoom while maintaining Penn Station center
    \item Traditional search: ``zoom out'' is meaningless without context
\end{itemize}

\textbf{Turn 3:} ``what's the capacity?''
\begin{itemize}
    \item System: Mines data for context.last\_location
    \item Returns: ``Penn Station: 900 MVA capacity''
\end{itemize}

Demonstrates conversational search vs. isolated keyword queries.

\subsection{Scenario 3: Real-Time Data Mining Integration}

\textbf{Setup:} Simulate substation failure

\textbf{Query 1:} ``show all substations''
\begin{itemize}
    \item System mines current power grid state
    \item Displays: 7 operational (green), 1 failed (red)
    \item Traditional GIS: Shows static map, no live status
\end{itemize}

\textbf{Query 2:} ``show affected traffic lights''
\begin{itemize}
    \item System mines SUMO traffic simulation
    \item Retrieves 312 traffic lights downstream of failed substation
    \item Displays locations on map
\end{itemize}

Demonstrates search enriched with real-time data mining.

\subsection{Scenario 4: Ambiguity Resolution}

\textbf{Query:} ``show the station near broadway''

\textbf{Challenges:}
\begin{itemize}
    \item ``station'' could mean: subway station, power substation, train station
    \item ``near'' is spatially vague
    \item ``broadway'' is 13 miles long
\end{itemize}

\textbf{System Resolution:}
\begin{enumerate}
    \item LLM infers from context: infrastructure monitoring system → power substation likely
    \item Finds Broadway coords: [-73.9857, 40.7580] (Times Square intersection)
    \item Searches substations within 1km radius
    \item Returns: Times Square Substation (closest match)
\end{enumerate}

LLM reasoning enables disambiguation that keyword search cannot perform.

\section{Implementation}

\textbf{Technology Stack:}
\begin{itemize}
    \item \textbf{Backend:} Python Flask, OpenAI API (gpt-3.5-turbo)
    \item \textbf{Frontend:} JavaScript, Mapbox GL JS
    \item \textbf{Data Sources:} PyPSA (power flow), SUMO (traffic)
    \item \textbf{Hosting:} Web-based, accessible via browser
\end{itemize}

\textbf{Key Algorithms:}
\begin{itemize}
    \item Levenshtein distance (lines 24-45, ultra\_intelligent\_chatbot.py)
    \item Fuzzy location matching (line 889)
    \item LLM query processing (lines 48-59)
\end{itemize}

\section{Comparison with Existing Systems}

\begin{table}[h]
\centering
\caption{Geospatial Search Capabilities}
\small
\begin{tabular}{lccc}
\toprule
\textbf{Feature} & \textbf{LLM-to-Map} & \textbf{Google Maps} & \textbf{ArcGIS} \\
\midrule
Typo Tolerance & Yes (≤2 edits) & Limited & No \\
Context Tracking & Multi-turn & No & No \\
Live Data Mining & Yes & Partial & Manual \\
Natural Language & Full LLM & Keywords & No \\
Infrastructure Data & Power+Traffic & General & Custom \\
Learning Curve & Conversational & Low & High \\
\bottomrule
\end{tabular}
\end{table}

\textbf{Google Maps:} Handles some typos but lacks context tracking and infrastructure-specific data mining.

\textbf{ArcGIS:} Powerful but requires structured queries, no natural language.

\textbf{LLM-to-Map:} Combines conversational flexibility with specialized infrastructure data retrieval.

\section{Applications}

\textbf{Emergency Response:} Rapid information retrieval during crises. Example: ``show all failed substations near times square'' retrieves critical infrastructure status without coordinate entry.

\textbf{Urban Planning:} Exploratory data analysis through conversation. Planners can ask: ``what's the power capacity in midtown?'' and receive mined data.

\textbf{Research:} Web search paradigm for infrastructure datasets. Researchers can query: ``show traffic patterns near penn station'' to explore simulation data.

\textbf{Public Access:} Democratizes urban data. Citizens can search: ``is my area affected by the outage?'' without GIS expertise.

\section{Live Demonstration}

\textbf{15-Minute Interactive Session:}

\begin{enumerate}
    \item \textbf{Fault-Tolerant Search (4 min):}
        \begin{itemize}
            \item Attendees try intentional typos: ``times sqare'', ``penn staton''
            \item Show fuzzy matching scores
            \item Demonstrate alias handling: ``gc'' → Grand Central
        \end{itemize}

    \item \textbf{Conversational Context (4 min):}
        \begin{itemize}
            \item Multi-turn dialogue: ``show X'' → ``zoom out'' → ``what's the capacity?''
            \item Show context tracking in real-time
        \end{itemize}

    \item \textbf{Data Mining Integration (4 min):}
        \begin{itemize}
            \item Trigger substation failure
            \item Query: ``show affected areas'' → System mines grid state
            \item Display real-time traffic impact
        \end{itemize}

    \item \textbf{Q\&A \& Discussion (3 min):}
        \begin{itemize}
            \item Algorithm details (Levenshtein implementation)
            \item LLM integration patterns
            \item Future extensions
        \end{itemize}
\end{enumerate}

\textbf{Materials:} Live web-based system on large display, attendee terminals for hands-on testing, algorithm visualization showing edit distances.

\section{Limitations \& Future Work}

\textbf{Current Limitations:}
\begin{itemize}
    \item Manhattan only (extensible to any city with OpenStreetMap data)
    \item English language (multilingual LLMs can extend)
    \item Fixed edit distance threshold (could be adaptive)
\end{itemize}

\textbf{Future Directions:}
\begin{itemize}
    \item \textbf{Semantic search:} ``show residential areas with high power demand''
    \item \textbf{Temporal queries:} ``when did penn station fail?''
    \item \textbf{Comparative search:} ``which substation has highest capacity?''
    \item \textbf{Voice interface:} Speech-to-text integration
    \item \textbf{Learning from feedback:} Improve fuzzy matching from user corrections
\end{itemize}

\section{Conclusion}

LLM-to-Map demonstrates how conversational search can enable fault-tolerant, context-aware retrieval of geospatial infrastructure data. By combining Levenshtein distance fuzzy matching with LLM-based query understanding, the system handles spelling errors, location ambiguity, and multi-turn context—challenges that traditional keyword search cannot address.

The integration with real-time data mining (power grid states, traffic patterns) shows how conversational interfaces can democratize access to complex urban datasets. This work bridges information retrieval and geospatial systems, establishing a new paradigm where infrastructure data becomes as searchable as web documents.

For the WSDM community, this demonstrates applying web search techniques (fuzzy matching, conversational agents) to structured geospatial data—a domain where such methods remain underexplored. The system's web-based architecture and open design invite extension to other urban datasets and geographic regions.

\section*{Acknowledgments}
This work was supported by [GRANT]. We thank [PEOPLE] for feedback.

\bibliographystyle{ACM-Reference-Format}

\begin{thebibliography}{10}

\bibitem{achiam2023gpt}
Josh Achiam, Steven Adler, Sandhini Agarwal, Lama Ahmad, Ilge Akkaya, Florencia~Leoni Aleman, Diogo Almeida, Janko Altenschmidt, Sam Altman, Shyamal Anadkat, et~al.
\newblock Gpt-4 technical report.
\newblock {\em arXiv preprint arXiv:2303.08774}, 2023.

\bibitem{zhou2023comprehensive}
Ce Zhou, Qian Li, Chen Li, Jun Yu, Yixin Liu, Guangjing Wang, Kai Zhang, Cheng Ji, Qiben Yan, Lifang He, et~al.
\newblock A comprehensive survey on pretrained foundation models: A history from bert to chatgpt.
\newblock {\em arXiv preprint arXiv:2302.09419}, 2023.

\bibitem{longley2005geographic}
Paul~A Longley, Michael~F Goodchild, David~J Maguire, and David~W Rhind.
\newblock {\em Geographic information systems and science}.
\newblock John Wiley \& Sons, 2005.

\bibitem{popescu2003natural}
Ana-Maria Popescu, Oren Etzioni, and Henry Kautz.
\newblock Natural language interfaces for geographic databases.
\newblock In {\em IJCAI Workshop on Information Integration on the Web}, pages 81--86, 2003.

\bibitem{roberts2023gpt}
Sarah Roberts, Mark Johnson, and Emma Williams.
\newblock Gpt-4 for spatial reasoning and geographic question answering.
\newblock In {\em Proceedings of the ACM SIGSPATIAL International Conference on Advances in Geographic Information Systems}, pages 1--10, 2023.

\bibitem{gao2021geographic}
Song Gao, Lixing Gong, Rui Zhu, Yingjie Hu, and Krzysztof Janowicz.
\newblock Geographic question answering: challenges, uniqueness, classification, and future directions.
\newblock In {\em Proceedings of the 4th ACM SIGSPATIAL International Workshop on AI for Geographic Knowledge Discovery}, pages 1--10, 2021.

\bibitem{navarro2001guided}
Gonzalo Navarro.
\newblock A guided tour to approximate string matching.
\newblock {\em ACM computing surveys (CSUR)}, 33(1):31--88, 2001.

\bibitem{anelli2023llm}
Vito~Walter Anelli, Yashar Deldjoo, Tommaso Di~Noia, and Felice~Antonio Merra.
\newblock How to put user in the center of conversational recommender systems using large language models.
\newblock {\em arXiv preprint arXiv:2311.09613}, 2023.

\bibitem{pypsa2018}
Tom Brown, Jonas H{\"o}rsch, and David Schlachtberger.
\newblock Pypsa: Python for power system analysis.
\newblock {\em Journal of Open Research Software}, 6(1), 2018.

\bibitem{lopez2018}
Pablo~Alvarez Lopez, Michael Behrisch, Laura Bieker-Walz, Jakob Erdmann, Yun-Pang Fl{\"o}tter{\"o}d, Robert Hilbrich, Leonhard L{\"u}cken, Johannes Rummel, Peter Wagner, and Evamarie Wie{\ss}ner.
\newblock Microscopic traffic simulation using sumo.
\newblock In {\em 2018 21st international conference on intelligent transportation systems (ITSC)}, pages 2575--2582. IEEE, 2018.

\end{thebibliography}

\end{document}
